% ============================================================================
\chapter{Introduction}
% ============================================================================

The \appname{} control application drives a \textbf{Delta~X~2} delta-arm robot used
for automated seeding of paper-pot trays. It is written in Rust with the Slint UI
toolkit and is designed for a dedicated operator terminal: a Raspberry~Pi~5 with a
7\,'' touch screen running in kiosk mode.

The application provides:

\begin{itemize}[itemsep=2pt]
  \item a guided, touch-friendly \textbf{seeding workflow} (select tray, load seeds,
        run, pause/resume, stop);
  \item a \textbf{calibration screen} for manual jogging of the X, Y and Z axes with
        selectable step sizes and mechanical homing;
  \item \textbf{software safety limits}: every movement is checked against configurable
        boundaries before any command is sent to the robot;
  \item live \textbf{position readout} and a permanent \textbf{status bar} with a
        connection indicator;
  \item plain-text configuration through a single \code{config.toml} file --- new tray
        layouts can be added without recompiling.
\end{itemize}

\section{Intended audience and structure}

Chapters~\ref{ch:installation} to~\ref{ch:configuration} are aimed at the
\emph{administrator} who installs and configures the system (keyboard and SSH access
assumed): Chapter~\ref{ch:installation} covers building the software,
Chapter~\ref{ch:rpi} its deployment on the Raspberry~Pi terminal, and
Chapter~\ref{ch:configuration} the \code{config.toml} settings.
Chapter~\ref{ch:usage} is aimed at the \emph{daily operator}, who interacts with the
machine exclusively through the touch screen.
Chapter~\ref{ch:troubleshooting} helps diagnose common problems.
Appendix~\ref{app:gcode} contains the complete G-code and communication protocol
specification of the Delta~X~2, and Appendix~\ref{app:limitations} lists current
development limitations.

\section{Typographic conventions}

\begin{itemize}[itemsep=2pt]
  \item \code{monospace} --- file names, commands, configuration keys, on-wire commands.
  \item \textbf{Bold} --- names of buttons and screens as they appear on the display.
  \item Shell commands prefixed with \code{\$} are typed in a terminal (locally or over SSH).
\end{itemize}

\begin{warning}
The Delta~X~2 is a moving machine. Keep hands and tools out of the working area while
the robot is powered. The software enforces motion limits, but software protections
are no substitute for physical caution. An additional hardware emergency-stop button
that cuts motor power is strongly recommended (see Appendix~\ref{app:limitations}).
\end{warning}
